\section{Discussion}
\label{sec:discussion}

\subsection{What we think the answer to the question is}

On the evidence so far: weight-tied looping does not pay for itself in a
zero-shot forecasting foundation model at this scale. It produces a gain of
about one and a half percent that barely clears a measured noise floor of
\NoiseFloorKFour\,\%, costs \CostRatio$\times$ the compute at training and
$k\times$ at inference, shows no monotone trend in $k$, and disappears when the
data pipeline is improved. Three interventions that add no depth each produce a
larger gain for a small fraction of the cost.

We want to be careful about the scope of that claim. It is a statement about
the \texttt{nano} configuration, this corpus, and this benchmark. It is not a
statement that weight tying is useless in general, and it is not a statement
about adaptive $k$ at inference, which we have not tested. What it does
support is a methodological point: in this setting, the architectural knob was
not where the headroom was, and the only way to find that out was to build the
cheap controls and measure them on the same footing.

\subsection{Why the negative result is the useful one}

Had we stopped after Phase~1 we would have reported ``$+2.1\,\%$ MASE from
looping at $k=4$'' from a single seed, which is both the luckier seed and an
effect smaller than the seed spread. The two things that changed the conclusion
cost almost nothing: training one extra seed, and building three baselines that
nobody asked for. Table~\ref{tab:levers} is now the bar that a future claim
about looping in this setting has to clear.

\subsection{Where the headroom appears to be}

The per-horizon decomposition points at the transition spectrum. Allowing
$\beta\in(0,2)$ --- a one-line change --- buys $+2.22\,\%$, and it is the only
lever whose benefit increases with horizon. \citet{ckda2026} shows that this is
a strict subset of the available move: with a real $\alpha \ge 0$ decay the
spectrum stays on the real line no matter what $\beta$ does, and only a
channel-wise \emph{signed} $\alpha$ combined with $\beta>1$ admits complex
eigenvalue pairs, which is what their periodic-waveform extrapolation results
turn on. Seasonality is periodic waveform continuation. The natural next
experiment is therefore not a larger $k$ but a signed channel-wise decay, and
it is cheap: a minor modification of kernels we already run, at
$96$--$97\,\%$ throughput.

\subsection{Threats to validity}

\paragraph{Reproduction gap.} Our baseline is \ReproGap\,\% behind the published
checkpoint \emph{in our harness}, on 15--20\,\% of the paper's compute. If the
gap is caused by something that interacts with depth --- corpus composition,
schedule length --- then conclusions about looping could change at the
published operating point. This is the single largest threat and we cannot
currently rule it out.

\paragraph{Harness scale.} Separately, our absolute MASE is a factor of
$1.419$ above the paper's for the \emph{same} checkpoint
(\S\ref{sec:background}). We argue this is an aggregation convention and show
that it cancels in every ratio we report, but the argument is algebra, not a
measurement; the measurement is pending. Until it is done, no absolute number
here should be read next to a published leaderboard.

\paragraph{Single seeds.} Only $k=4$ has two seeds. The lever arms are
single-seed, so the $1.70$--$3.41\,\%$ figures are larger than the $k=1$ seed
spread ($0.77\,\%$) but have not been replicated. \texttt{AUGNEG} is being run
with two seeds for this reason.

\paragraph{One benchmark, one model size.} All results are GIFT-Eval
\texttt{nano}. No claim is made about scaling.

\paragraph{Mechanism inference from three arms.} The horizon-profile argument
in \S\ref{sec:results} compares three single arms. It is suggestive and it was
pre-registered, but it is not a controlled test of mechanism.
